% Full independent literal audits by root and B passed; source note retained.
\subsection{Large corners with four objects and seventeen arrows}

\begin{lemma}[Fourth homology at excess three]
\label{C:arrow-excess-three-H4}
If a connected finite category has $k$ objects and at most
$2k+2$ arrows, then its fourth rational homology is zero.
\end{lemma}
\begin{proof}
Delete retract objects. Each deletion removes at least three
arrows, preserves homotopy type, and preserves the stated bound.
In the resulting category, a proper strong component on $t\geq2$
objects has at least $t^2+t$ internal arrows, since all its
identities lie outside its minimal ideal. Connectivity outside
that component gives a total of at least
$2k+t^2-1\geq2k+3$, a contradiction. If the whole category is
strongly connected, it has either one object and at most four
arrows, or two objects and exactly six arrows. The first case
is rationally acyclic by the small-monoid classification.
In the second case the minimal band has one row and one column
at each object and trivial maximal group. The minimal arrows
from a chosen object give a natural contraction, as in the
proof of Lemma~\ref{C:one-row-source-block-collapse}.

It remains to consider acyclic support. In the notation used
for the excess bounds, $D+P+R\leq3$. If the diagonal algebra
is semisimple, a strict chain has at most four objects, because
a five-object chain has support cycle rank six. The relative
bar therefore has length at most three and $H_4=0$.
Otherwise $D\geq2$. Every corner has order at most four and
is rationally acyclic. If $P=0$, the internal bars for a strict
chain have rational homology only in degree zero. Thus the
actual support functor is a rational homology equivalence;
the support has no four-simplex and its fourth homology is zero.
The sole remaining possibility is $D=2$, $P=1$, $R=0$.
There is one corner of order three, all other corners are
trivial, and the support is a tree with one doubled hom-set.
The trivial modules of every order-three monoid have rational
projective dimension at most one, by the explicit order-three
algebra list in Proposition~\ref{C:semisimple-ends-small-middle}:
the identity-adjoined two-element one-sided zero bands have
a projective one-term augmentation kernel on the nonprojective
side, and the other cases have projective trivial modules.
Consequently the only
possibly nontrivial edge collage, an ordinary two-state act
cone at that corner, has zero homology in degrees at least
three. The other edge collages are graphs or singleton joins.
Their corner intersections are rationally acyclic, so
Mayer--Vietoris gives zero fourth homology for their union.
\end{proof}

\begin{lemma}[A singleton arrow with a small outside budget]
\label{C:singleton-extreme-three-excess}
Let $C$ be connected and have a sink object $v$. Write $M$
for its corner, $D=C\setminus\{v\}$, and
$\ell=|\operatorname{Ob}D|\geq1$. Suppose that some incoming
hom-set at $v$ is a singleton and that
\[
 |D|+\sum_{x\in D}|C(x,v)|\leq2\ell+3.
\]
Then all rational degree-two products on $BC$ vanish.
The opposite assertion holds for a source object.
\end{lemma}
\begin{proof}
If $H^2(BD;\mathbb Q)=0$, the singleton fixed-arrow lemma
applies directly. Suppose instead this group is nonzero,
and let $D_1,\ldots,D_c$ be the connected components of $D$.
At least one component on $t$ objects has at least $2t+2$
arrows by Lemma~\ref{C:arrow-excess-two-H2}; each other
component has at least $2t-1$ arrows. Hence
\[
 |D|\geq2\ell-c+3.
\]
Connectivity requires at least one incoming arrow from each
component. The assumed budget forces equality throughout:
there is precisely one arrow from each component to $v$,
and the sum of the component excesses is exactly three.
Each component therefore satisfies the preceding lemma and
has zero fourth rational homology.

Let $w_i$ be the unique object of $D_i$ with an arrow to $v$.
It has no incoming cross arrow in $D_i$, since composing one
with $w_i\to v$ would give an additional arrow into $v$.
Its corner has order at most four: connectivity in $D_i$
and excess at most three leave at most three extra corner
arrows. In particular that corner is rationally acyclic.
Let $Z$ be the full star on $v,w_1,\ldots,w_c$. Its legs are
singletons. Each branch is the join of $BM$ with a rationally
acyclic connected corner space, so is rationally acyclic.
The branches meet along $BM$. Their Mayer--Vietoris sequence
gives $H^2(BZ;\mathbb Q)=0$, since $H^1(BM;\mathbb Q)=0$.

The nerve of $C$ is the actual union of this star and the
component nerves $ND_i$, meeting at their corner nerves
$N\operatorname{End}(w_i)$. Indeed, a chain using $v$ can
use no other object before $w_i$. These intersections are
rationally acyclic. Fourth-cohomology restriction to all
these pieces is injective. On $BZ$ a degree-two class is
zero, and on each $BD_i$ its square is zero because $H^4=0$.
The global square, and hence all mixed degree-two products
by polarization, is zero. Opposition proves the dual claim.
\end{proof}

\begin{proposition}\label{C:seventeen-four-large-corner}
A four-object category with at most seventeen arrows and an
endomorphism monoid of order at least eight cannot have type
$\mathbb CP^2$.
\end{proposition}
\begin{proof}
Assume it has that type and is connected. A retract deletion
leaves at most fourteen arrows, already excluded. Thus no
object is a retract of another. Let $v$ have corner $M$ of
order $q\geq8$.

First exclude a nontrivial strong component. One not containing
$v$ has at least six internal arrows. The other two corners
and the necessary interfaces give at least
$6+8+1+3=18$ arrows. A strong component of at least three
objects containing $v$ requires at least
$8+2+2+6+1+3=22$ arrows if proper, and still more if it is
the whole category. A strong pair containing $v$ has at least
$q+4$ internal arrows, and its two-object complement and
interfaces require at least five more. The only possible
case is $q=8$, internal size twelve, and total size seventeen.
It is extreme: a nonextreme pair needs two expanded interfaces
and their composite, giving at least seven outside arrows.
Its other corner is $\{1,e\}$, both internal cross hom-sets
are singletons, and their composite at that corner is $e$.
If the pair is a source, all incoming arrows at the small
corner are internal and fixed by $e$; if a sink, use outgoing
arrows. The one-sided corner-deletion lemma applies, leaving
a smaller forbidden category. We may assume acyclic support.

Let $d$ be the sum of the corner excesses away from $v$.
With $P,R$ denoting parallel excess and support cycle rank,
the arrow budget is
\[
 d+P+R\leq11-q\leq3.
\]
If $P=0$, every degree-two class pulls back from the support
poset, whose fourth homology is zero. We may assume $P\geq1$.

\paragraph{A nonextreme large corner.}
Suppose first that $v$ has both entrances and exits.
Then $R\geq1$, and hence $d\leq1$. If $d=1$, necessarily
$q=8$ and $P=R=1$. The support is one triangle through
$v$ and one leaf edge. The large corner is the middle of
the triangle. Its fourth object can attach only at one of
the two extremes and away from that extreme: any attachment
through the middle or toward the triangle would force another
comparison and a second cycle. Therefore the single order-two
corner, wherever it lies, is extreme in the whole category.
Its exact orbit quotient makes it trivial, preserving rational
cohomology and not increasing any hom-set. There is now at
most one doubled hom-set. One of the two adjacent hom-sets
at $v$ is a singleton, and its neighboring corner is trivial.
Singleton cofinal deletion removes $v$.
The remaining three corners are trivial, so their acyclic
category nerve has dimension at most two and zero fourth
homology. This contradicts the preserved rational square.

We may thus assume all other corners are trivial. If the
total number of incoming arrows at $v$ is one, or the total
number of outgoing arrows is one, the same cofinal deletion
applies, again leaving a nerve of dimension at most two.
Assume both totals are at least two. There are at
most six cross arrows in all.
If $v$ has only two neighbors, they lie on opposite sides.
The two incident hom-sets each have size two, their composite
hom-set is singleton, and the last object attaches by one
singleton edge. This uses all six arrows. Transitivity puts
that leaf edge at a trivial extreme, so the full triangle and
leaf edge form an actual nerve union along a point. The triangle
has a trivial extreme whose two stalks have sizes one and two;
Lemma~\ref{C:two-states-objectwise-act-cups} kills its products.
The leaf edge contributes no higher cohomology.
If $v$ has three neighbors, two are on one side and the third
on the other. Counting the two forced composite hom-sets
again uses exactly six arrows: the two-side hom-sets are
singletons, the other incident hom-set has size two, and
the composites are singletons. At the opposite trivial
extreme all stalks have size at most two. The same categorical
act-cone theorem kills all products. These exhaust the
nonextreme case.

\paragraph{An extreme large corner.}
Take opposites to make $v$ a sink. If it has a singleton
incoming hom-set, its outside budget is at most
$17-q\leq9=2\cdot3+3$, so the preceding lemma applies.
Assume henceforth every incoming hom-set has size at least two.
Let $r$ be the number of its neighbors.

For $r=3$, the entire budget forces $q=8$, all other corners
trivial, and three legs of size two. Identifying the trivial
leaves gives an ordinary act cone over $M$, adding only circles.
It inherits bounded rational homology and second-cohomology
rank one, contradicting Lemma~\ref{C:nine-bounded-rank-one}.

For $r=1$, denote its sole neighbor by $a$ and put
$s=|C(a,v)|$. The complement $D$ is connected, and no
cross arrow in $D$ enters $a$, by transitivity. The nerves
of $D$ and the two-object collage $Y$ on $a,v$ cover $C$
and meet in $B\operatorname{End}(a)$.
If $s=2$, then $|D|\leq7$ and its second rational cohomology
is zero by Lemma~\ref{C:three-small-H2-zero}. The shared
corner has order at most three and is rationally acyclic.
The universal two-cross-arrow theorem makes every degree-two
product on $BY$ zero. Fourth-cohomology restriction to $D,Y$
is injective, proving the contradiction.
If $s\geq3$, connectivity gives $|D|\geq5$, so $s\leq4$
and $q\leq9$. All corners of $D$ have order at most two.
Also $|D|\leq6$, so its second cohomology is zero by the
same three-object lemma. Its semisimple relative bar has
no high homology, and the
shared corner is rationally acyclic. Thus $Y$ inherits
bounded rational homology and second-cohomology rank one.
Exact coinvariants at its small corner give an ordinary
act over $M$ with at most four states. The same bounded
rank-one lemma excludes it.

It remains that $r=2$. Write $a,b$ for its neighbors and
$w$ for the fourth object. A cross arrow involving $w$
must point from $a$ or $b$ to $w$; an arrow in the opposite
direction would force $w\to v$. The two incoming hom-sets,
the remaining connection, and the corner identities force
$q\leq9$. If $q=9$ there is no remaining excess. If $q=8$
there is at most one further corner or cross arrow beyond
these minimum contributions.

If there is just one nonempty hom-set away from $v$, then
the category is the actual union of its central two-leg star
and the edge to $w$, along the corner at $a$ or $b$.
An order-two corner is the only possible extra corner, and
their common corner is rationally acyclic. A singleton edge
is a rationally acyclic join. If that edge has size two,
both endpoints are trivial and it contributes a circle,
contradicting simple connectivity. In the singleton case
the central star inherits bounded rational homology and
rank-one second cohomology. Its leaves have order at most
two, so exact leaf coinvariants followed by point identification
give an ordinary act over $M$, again excluded by the bounded
rank-one lemma.

Otherwise there are exactly two hom-sets away from $v$,
both singletons, $q=8$, all other corners are trivial,
and both incoming hom-sets at $v$ have size two.
If these two hom-sets are $a\to w$ and $b\to w$, the
support is an induced four-cycle. The support functor is
onto on fundamental groups, contradicting simple connectivity.
Otherwise, after swapping $a,b$, they are $a\to b$ and
$a\to w$. Delete the trivial leaf $w$. The remaining
triangle has hom-set sizes $(1,2,2)$, and composition is
an equivariant map between two-element sets. If constant,
its value is a fixed incoming arrow at $M$, so the fixed-arrow
lemma, Lemma~\ref{C:fixed-incoming-cups}, applies with $D$
the trivial two-object interval $a\to b$ and that constant
image as the $M$-fixed arrow into $v$. If it is
bijective, the actual inverse composition map deletes $b$
toward $a$, leaving an ordinary two-state act cone.
Both alternatives have zero degree-two products.
This exhausts all budgets and proves the proposition.
\end{proof}
